\subsection{The local five-dimensional quotient on three rows}

\begin{lemma}[An affine pencil cannot give a lone positive even Tor group]
\label{C:drifting-pencil-Tor}
Let $A$ be the rational algebra with basis $1,x,y,z,c$, with
\[
 x^2=z,\qquad xz=zx=y^2=c,\qquad
 xy=\epsilon c,\qquad yx=(1-\epsilon)c,
 \qquad \epsilon\in\{0,1\},
\]
and all other products between $x,y,z,c$ zero.
For $a\in\mathbb Q$, let $V_a$ be the two-dimensional module
on which $y$ acts as a nonzero square-zero Jordan block, $x$
acts as $a$ times that block, and $z,c$ act as zero.
There is no finite-dimensional right $A$-module $W$ with
\[
 \operatorname{Tor}_1^A(W,V_a)=0,\qquad
 \dim\operatorname{Tor}_2^A(W,V_a)=1,\qquad
 \operatorname{Tor}_n^A(W,V_a)=0\quad(n\geq3).
\]
\end{lemma}
\begin{proof}
Put $\eta=1-\epsilon$, $s=\epsilon-\eta$, and
$t_n=a+ns$.  A free resolution of $V_a$ has one copy of $A$
in every degree and differentials, specified by the images of $1$,
\[
 d_{2n+1}(1)=x-t_ny,\qquad
 d_{2n+2}(1)=y+(t_n-\eta)z\qquad(n\geq0).
\]
These are maps of left modules, hence right multiplications by
the displayed elements.  Direct multiplication gives
\begin{align*}
 (y+(t-\eta)z)(x-ty)&=0,\\
 (x-(t+s)y)(y+(t-\eta)z)&=0.
\end{align*}
The left ideal generated by $x-ty$ has basis
$x-ty,z,c$ for every rational $t$, and the left ideal generated
by $y+(t-\eta)z$ has basis $y+(t-\eta)z,c$.
Thus successive images have dimensions three and two, their
composites are zero, and they are exact in the five-dimensional
free terms.  The initial cokernel is $V_a$, so this is the
claimed complete resolution.

After tensoring with $W$, write $d=\dim W$ and $r_i$ for the
rank of the $i$th differential.  The asserted Tor dimensions
give
\[
 r_1+r_2=d,\qquad r_2+r_3=d-1,
 \qquad r_n+r_{n+1}=d\quad(n\geq3).
\]
In particular $r_1-r_3=1$, while
$r_3=r_5=r_7=\cdots$.
The odd differentials are the matrix pencil
$R_x-t_nR_y$, where $R_x,R_y$ are the right-action matrices
on $W$.  Since $s=\pm1$, these parameters are distinct.
A rational matrix pencil has its maximal rank at all but
finitely many rational parameters: choose a nonzero minor of
maximal size and exclude its finitely many roots.
Consequently $r_3$ equals that maximal rank.  Its specialization
at $t_0=a$ has rank $r_1\leq r_3$, contradicting $r_1-r_3=1$.
The argument uses the stated vanishing in every degree at least
three, not just an unspecified eventual tail with additional
intervening homology.
\end{proof}

\begin{proposition}[Every local five-element contracted basis on three rows]
\label{C:local-five-three-row-Tor}
Let $A$ be a local contracted rational algebra of a finite
monoid, with $\dim A=5$.  Suppose $V$ is a two-dimensional
$A$-module with a fixed augmentation basis of a three-point
set, in which every distinguished nonzero monoid basis element
acts by the augmentation matrix of a transformation of that
set.  The transformations themselves need not descend through
the contracted zero; only their augmentation operators are
required to give the stated $A$-module.
There is no finite-dimensional right module $W$
with
\[
 \operatorname{Tor}_1^A(W,V)=0,\qquad
 \dim\operatorname{Tor}_2^A(W,V)=1,\qquad
 \operatorname{Tor}_n^A(W,V)=0\quad(n\geq3).
\]
\end{proposition}
\begin{proof}
We give the exact finite basis calculation and its independent
verification, including the only infinite-resolution issue.
In a local contracted monoid algebra there is no nonidentity
idempotent basis element and no nontrivial finite subgroup:
either would give a nontrivial idempotent in the algebra.
Every nonidentity nonzero monoid element therefore has zero
idempotent power.  The span of these four elements is a
nilpotent ideal and is the radical $J$, with quotient
$A/J=\mathbb Q$.  Its nilpotence is also checked literally
in the finite enumeration below.

If $V$ is semisimple, it is the sum of two copies of the simple
module and its degree-two Tor has even dimension.  Otherwise
$JV$ is one-dimensional and $J^2V=0$.  All radical operators
are scalar multiples of one nonzero square-zero matrix $N$.
A three-point transformation with nonzero nilpotent augmentation
operator is a length-two chain; relabeling makes it $(0,0,1)$.
In its augmentation basis $N=\left(\begin{smallmatrix}0&1\\0&0\end{smallmatrix}\right)$.
The actual three-point transformations whose augmentation is
a multiple of $N$ are the three constants, $(0,0,1)$, and
$(1,1,0)$, with coefficients $0,1,-1$ respectively.
Thus every nonidentity distinguished basis element acts by
$0,N$, or $-N$, and its coefficient vanishes on $J^2$.

Start from all $1373$ isomorphism-and-opposite classes in the
complete order-six census
\path{code/outputs/c_monoids_6.tsv}, add literal opposites,
retain tables with a zero, and contract it.  Computing powers
of the four-dimensional nonunit ideal leaves $156$ distinct
literal local bases.  On each basis enumerate all nonzero
assignments of coefficients in $\{-1,0,1\}$ to its four
nonidentity elements and impose the multiplication equations
on their two-by-two matrices.  There are $3600$ assignments.
Relabeling the four nonidentity basis elements and allowing the
common change $N\mapsto-N$ leaves $1286$ pairs $(A,V)$.
This deliberately retains assignments that might not admit
an affine lift to a literal three-point action; overenumeration
cannot discard a possible source.

Exact rational projective covers and explicit kernel bases give
the following complete partition:
\[
\begin{array}{l|r}
 \text{verified obstruction}&\text{pairs}\\\hline
 \text{a simple direct summand in a syzygy}&414\\
 \text{a syzygy with top dimension greater than radical dimension}&859\\
 \text{a split recurrence of the first syzygy}&7\\
 \text{the two affine-pencil algebras of Lemma~\ref{C:drifting-pencil-Tor}}&6.
\end{array}
\]
The first line supplies actual inclusion and retraction matrices.
If the proposed Tor profile held, $W$ would be nonprojective;
its Tor with the local simple module is then nonzero in every
degree, by Lemma~\ref{B:local-minimal-Tor}.  Dimension shifting
therefore excludes that line.  The second line is excluded by
Lemma~\ref{B:top-heavy-Tor}, again using dimension shifting.
The seven recurrences give actual split maps from the first
syzygy into a later syzygy.  They are excluded by
Lemma~\ref{B:split-syzygy-persistence}.

For the last six records, a literal permutation of the radical
basis gives exactly the multiplication in
Lemma~\ref{C:drifting-pencil-Tor}; the row module has
$a\in\{-1,0,1\}$.  That lemma supplies the complete resolution
and the rank obstruction.  Equality of successive module
dimensions is not treated as module periodicity: in these six
records the pencil parameter actually changes.

The constructor is \path{code/a_five_local_all_rows.py}.
The separate checker is
\begin{quote}\small
\path{code/a_verify_five_local_all_rows.py}.
\end{quote}
It imports none of the search, cover, or syzygy implementations
and rebuilds the
entire $3600$-assignment source coverage directly from the
census, using nilpotence of the full ideal in place of the
constructor's individual power test, and independently imposes
all matrix multiplication equations.  It checks every one of
the $1372$ stored projective-prefix steps by surjectivity,
equivariance, zero composite, and complementary exact ranks;
it checks every inclusion and retraction and every top dimension.
For the six exceptional records it verifies the two zero
composites as polynomial identities and exhibits constant unit
minors proving ranks three and two for every parameter.
The result
\path{code/outputs/a_five_local_all_rows_independently_verified.json}
is \texttt{PASS} with precisely the displayed counts and no
unresolved pair.  The source census and certificate hashes
are recorded there, so the finite scope is explicit.
\end{proof}

\begin{corollary}[Local quotients with five outside elements]
\label{C:local-five-two-column-act-cones}
Let $M$ have rectangular-band minimal ideal of shape $3\times2$,
with exactly five outside elements, and suppose
$\mathbb QM/\mathbb QK$ is local.
For every nonempty finite right $M$-set $S$,
$B\mathcal C(M,S)$ cannot have the rational homology and
cohomology ring of $\mathbb CP^2$.
In particular it cannot be homotopy equivalent to $\mathbb CP^2$.
\end{corollary}
\begin{proof}
A $K$-fixed state makes every degree-two square zero, as before.
Otherwise Lemma~\ref{C:two-column-projective-orbit-kernel} and
constant-action compression give a monoid $N$ with the same
homology as the act cone in every degree at least three.
Its contracted algebra and three-point row augmentation module
are the original $A=\mathbb QM/\mathbb QK$ and $V$.
The band--Tor formula would give
$\operatorname{Tor}_1^A(W,V)=0$,
$\dim\operatorname{Tor}_2^A(W,V)=1$, and zero Tor in every
higher degree, for the new column augmentation module $W$.
Proposition~\ref{C:local-five-three-row-Tor} excludes this.
Only the rational homology profile was used in this branch.
\end{proof}
